% =============================================================================
% tgp_nbody_results_clean.tex
% Clean publication-ready draft: novel results from TGP N-body analysis
%
% Self-contained document (standalone compilation with pdflatex).
% Can also be \input{}-ed into the main TGP manuscript as a section/appendix.
%
% Notation synchronized with: nbody/pairwise.py, nbody/eom_tgp.py,
% nbody/three_body_force_exact.py, nbody/tgp_nbody_lagrangian_eom.tex.
%
% Mateusz Serafin, 2026
% =============================================================================

\documentclass[11pt,a4paper]{article}

\usepackage[utf8]{inputenc}
\usepackage[T1]{fontenc}
\usepackage{amsmath,amssymb,amsthm}
\usepackage{mathtools}
\usepackage{booktabs}
\usepackage[pdfusetitle,hidelinks]{hyperref}
\usepackage[margin=2.5cm]{geometry}
\usepackage{natbib}
\bibliographystyle{unsrtnat}
\pdfstringdefDisableCommands{%
  \def\textup#1{#1}%
  \def\textsc#1{#1}%
  \def\msp{m\_sp}%
  \def\Ceff{C\_eff}%
}
\usepackage{enumitem}

\newcommand{\vx}{\mathbf{x}}
\newcommand{\vF}{\mathbf{F}}
\newcommand{\vv}{\mathbf{v}}
\newcommand{\vL}{\mathbf{L}}
\newcommand{\vP}{\mathbf{P}}
\newcommand{\bR}{\mathbb{R}}
\newcommand{\dij}{d_{ij}}
\newcommand{\msp}{m_{\rm sp}}
\newcommand{\lmax}{\lambda_{\max}}
\newcommand{\drep}{d_{\rm rep}}
\newcommand{\dwell}{d_{\rm well}}
\newcommand{\Ceff}{C_{\rm eff}}

\newtheorem{result}{Result}
\newtheorem{prediction}{Prediction}

\title{%
  \textbf{N-body dynamics of Topology-Generated Potential:}\\[4pt]
  Earnshaw violation, chaos suppression,\\
  and the EFT origin of Yukawa coupling%
}
\author{Mateusz Serafin}
\date{April 2026}

\begin{document}
\maketitle

\begin{abstract}
We derive and verify numerically the complete $N$-body dynamics arising
from the TGP scalar-field Lagrangian $\mathcal{L}[\Phi]$ with quartic
self-interaction.  The potential decomposes as $V = V_2 + V_3^{\rm irr}$,
where $V_2$ is a closed-form pairwise interaction (gradient + repulsive +
confining) and $V_3$ is an irreducible three-body term computed via a
Feynman two-parameter integral.

Seven principal results emerge.  (1)~TGP breaks the Earnshaw theorem:
equilateral three-body configurations admit stable static equilibria with
all normal modes positive ($\omega^2 > 0$), whereas Newton gives marginal
modes ($\omega^2 \approx 0$) for all tested parameter pairs $(\beta,C)$.
(2)~TGP suppresses three-body chaos by 35--45\% (measured via the
maximal Lyapunov exponent) for $\beta \gtrsim 0.025$ at initial
conditions involving close encounters; the effect is IC-dependent and
enhanced by $V_3$.  (3)~The Yukawa source coupling $\Ceff$ is derived
from EFT projection of the classical topological defect.
(4)~Classical soliton interaction is \emph{not} Yukawa (oscillatory,
power-law decay), establishing that the Yukawa law requires the EFT mass
gap.  (5)~The irreducible three-body integral admits a semi-analytic
multipole representation with controlled accuracy ($<0.01\%$ at
$L_{\max}=15$).  (6)~The ratio $|V_3/V_2|$ is mapped across the
dimensionless parameter $t = \msp\,d$, showing perturbative behaviour
($\sim1\%$) in the Lyapunov regime and exponential suppression for
$t > 3$.  (7)~Closed-form expressions for the equilibrium distances
$\drep$ and $\dwell$, the radial normal-mode frequency
$\omega_{\rm rad}^2$, and the TGP Hill sphere radius are derived and
verified; TGP generically expands the Hill sphere relative to Newton.
(8)~For unequal masses, the phase diagram is governed by a critical
coupling $\beta_{\rm crit} = \tfrac{3}{2}\sqrt{\gamma(C_1+C_2)}$;
the escape velocity and a three-body breathing mode are derived in
closed form.
\end{abstract}

\tableofcontents

% =====================================================================
\section{Introduction}
\label{sec:introduction}
% =====================================================================

The gravitational $N$-body problem has been a cornerstone of
mathematical physics since Newton.  In the restricted three-body case,
the problem exhibits chaotic trajectories, non-integrability (Bruns
theorem~\citep{Bruns1887}), and the impossibility of stable static
equilibria (Earnshaw's theorem~\citep{Earnshaw1842}).  Modern
treatments range from the classical monographs of
\citet{Valtonen2006} and \citet{Heggie2003} to statistical approaches
for chaotic outcomes~\citep{StoneLeigh2019,Manwadkar2021}.  The Burrau
problem~\citep{Burrau1913,Szebehely1967} remains a standard benchmark
for numerical integrators.

In modified-gravity theories, $N$-body dynamics are altered by
additional scalar or vector degrees of freedom.  Cosmological $N$-body
codes incorporating scalar-field modifications include
MG-GADGET~\citep{Puchwein2013}, ECOSMOG~\citep{Li2012}, and
ISIS~\citep{Llinares2014}; a systematic comparison was carried out
by~\citet{Winther2015}.  These codes focus on large-scale structure
($\gtrsim$~Mpc), and screening mechanisms
(chameleon~\citep{Khoury2004}, Vainshtein~\citep{Vainshtein1972})
play a central role.  At the fundamental level, scalar-tensor
theories~\citep{Damour1992} introduce a Yukawa correction to the
Newtonian potential, constrained by fifth-force
experiments~\citep{Adelberger2003,Murata2015} and solar-system
tests~\citep{Will2014,Bertotti2003}.

A separate line of research concerns \emph{few-body} systems with
short-range (Yukawa-type) interactions, where three-body forces can
produce qualitatively new phenomena.  The Efimov
effect~\citep{Efimov1970} --- an infinite tower of three-body bound
states near a two-body resonance --- illustrates how irreducible
three-body physics can dominate; see~\citet{Braaten2006}
and~\citet{Naidon2017} for reviews.  While the Efimov effect is
typically studied in quantum mechanics, the underlying principle ---
that non-additive three-body forces qualitatively alter the spectrum
--- is relevant to any system with short-range interactions.

The Topology-Generated Potential (TGP)~\citep{Serafin2026} is a
scalar-field theory in which space is emergent from a discrete
substrate.  The field Lagrangian naturally produces a quartic
self-interaction that, upon expansion around solitonic sources, yields:
(i)~a pairwise potential $V_2$ combining Newtonian attraction, a
repulsive barrier, and short-range confinement;  (ii)~an irreducible
three-body potential $V_3$ arising from the $\Phi^4$ vertex.  This
paper derives the complete $N$-body dynamics from the TGP Lagrangian
and presents eight principal results, each verified by a dedicated
regression suite (59/59 PASS, 102 active scripts).

The paper is organized as follows.
Section~\ref{sec:eom} derives the equations of motion.
Section~\ref{sec:earnshaw} establishes Earnshaw violation (Result~1).
Section~\ref{sec:chaos} quantifies chaos suppression (Result~2).
Section~\ref{sec:pathC} derives the EFT Yukawa coupling (Result~3).
Section~\ref{sec:soliton-interaction} contrasts classical soliton
interaction with the Yukawa law (Result~4).
Sections~\ref{sec:multipole}--\ref{sec:v3v2} analyze the three-body
integral and its regime map (Results~5--6).
Section~\ref{sec:analytical} derives analytical equilibria and the
Hill criterion (Result~7).
Section~\ref{sec:unequal} treats unequal masses and the phase diagram
(Result~8).
Section~\ref{sec:summary} summarizes the deductive chain.

% =====================================================================
\section{Equations of motion from TGP Lagrangian}
\label{sec:eom}
% =====================================================================

\subsection{Setup}

We consider $N$ point sources with positions $\vx_i \in \bR^3$ and
source strengths $C_i > 0$ serving as both gravitational ``charge'' and
inertial mass ($m_i \equiv C_i$).  The TGP Lagrangian yields a
total potential $V = V_2 + V_3$ with the following closed-form expressions.

\subsection{Pairwise potential $V_2$}

For bodies $i \neq j$ at separation $\dij = |\vx_i - \vx_j|$,
\begin{equation}
  \boxed{
  V_{2,ij}(\dij)
  = \underbrace{-\frac{4\pi\,C_i C_j}{\dij}}_{\text{gradient (attractive)}}
  + \underbrace{\frac{8\pi\,\beta\,C_i C_j}{\dij^2}}_{\text{repulsive ($\beta$-term)}}
  - \underbrace{\frac{12\pi\,\gamma\,C_i C_j\,(C_i + C_j)}{\dij^3}}_{\text{confining ($\gamma$-term)}},
  }
  \label{eq:V2}
\end{equation}
with total $V_2 = \sum_{i<j} V_{2,ij}$.  The three terms arise from the
$(\nabla\Phi)^2$, $\Phi^3$, and $\Phi^4$ sectors of the TGP action,
respectively.  The hierarchy $1/d \gg 1/d^2 \gg 1/d^3$ at large
separation ensures Newtonian behaviour at long range; at short distances
the $\beta$-repulsion creates a barrier $\drep$ (when it exists) and the
$\gamma$-confinement creates a well $\dwell$.

\subsection{Irreducible three-body potential $V_3$}

For each unordered triplet $\{i,j,k\}$, the irreducible three-body sector
receives contributions from both $\Phi^3$ and $\Phi^4$:
\begin{equation}
  \boxed{
  V_{3,ijk}
  = (2\beta - 6\gamma)\,C_i C_j C_k \; I_Y(d_{ij}, d_{ik}, d_{jk};\, \msp),
  }
  \label{eq:V3}
\end{equation}
where $I_Y$ is the triple-overlap integral in the Feynman
(Schwinger-parameter) representation:
\begin{equation}
  I_Y = 2 \int_{\Delta_2}
    \Delta^{-3/2}\;
    K_0\!\Bigl(\msp\sqrt{Q/\Delta}\Bigr)\;
    \mathrm{d}\alpha,
  \label{eq:IY-feynman}
\end{equation}
with $\Delta = \alpha_1\alpha_2 + \alpha_1\alpha_3 + \alpha_2\alpha_3$,
$Q = \alpha_2 d_{ij}^2 + \alpha_1 d_{ik}^2 + \alpha_3 d_{jk}^2$,
the integral being over the standard 2-simplex
$\Delta_2 = \{\alpha \geq 0,\; \alpha_1+\alpha_2+\alpha_3=1\}$.
The screening mass is
\begin{equation}
  \msp = \sqrt{3\gamma - 2\beta}.
  \label{eq:msp}
\end{equation}
In the unscreened (Coulomb) limit $\msp \to 0$:
$I \to 8\pi^2/(d_{ij}+d_{ik}+d_{jk})$.

\subsection{Equations of motion}

With $V = V_2 + \sum_{i<j<k} V_{3,ijk}$, the Euler--Lagrange equations
give
\begin{equation}
  C_i\,\ddot\vx_i = \vF_i = -\nabla_{\vx_i} V,
  \label{eq:eom}
\end{equation}
with exact translational invariance ($\sum_i \vF_i = 0$),
$SO(3)$ rotational invariance ($\sum_i \vx_i \times \vF_i = 0$), and
Hamiltonian structure ($H = T + V$ conserved).

\medskip
\noindent\textbf{Remark.}  $V_2$ and $V_3$ are the \emph{only} terms arising
from the standard TGP $\Phi^4$ Lagrangian for point-source superposition;
there is no free structural parameter beyond $(\beta, \gamma, \{C_i\})$.

% =====================================================================
\section{Result 1: Breaking Earnshaw's theorem}
\label{sec:earnshaw}
% =====================================================================

\begin{result}[Earnshaw violation]
For equilateral three-body configurations at the confining equilibrium
$\dwell$, TGP produces \textbf{stable} static equilibria (all
$\omega^2 > 0$), while Newton yields marginal modes
($\omega^2 \approx 0$) for all tested $(\beta, C)$ pairs.
\end{result}

\subsection{Background}

Earnshaw's theorem~\citep{Earnshaw1842} states that no static
equilibrium of point masses under an inverse-square-law force can be
stable in all directions.
In TGP, the higher-order terms ($1/d^2$ repulsion, $1/d^3$ confinement)
create an effective well $\dwell$ where $V'_2(\dwell) = 0$ with
$V''_2(\dwell) > 0$.  The question is whether the full $3N$-dimensional
Hessian $\partial^2 V/\partial x_{i\alpha}\partial x_{j\beta}$ has all
non-translational eigenvalues positive.

\subsection{Normal-mode analysis}

We compute the mass-weighted dynamical matrix
$D = M^{-1/2}\,H\,M^{-1/2}$ where $H$ is the $3N \times 3N$ Hessian
and $M = \mathrm{diag}(C_1, C_1, C_1, C_2, \ldots)$.  Normal-mode
frequencies satisfy $\omega_k^2 = \text{eigenvalues of } D$, with
translational zero modes removed.

\subsection{Results}

For an equilateral triangle ($N=3$) on $\dwell$, scanning
$\beta \in [0.5, 5.0]$ and $C \in [0.05, 0.50]$:

\begin{center}
\begin{tabular}{lccl}
  \toprule
  \textbf{Mode type} & \textbf{Newton $\omega^2$} & \textbf{TGP $\omega^2$}
    & \textbf{Character} \\
  \midrule
  Breathing (radial) & $\sim 0$ (marginal) & $0.0083$ & Radial-fraction $= 1.0$ \\
  Mixed $\times 2$ ($C_3$ degen.)  & $\sim 0$ & $0.0042$ & Tangential + radial \\
  \bottomrule
\end{tabular}
\end{center}

\begin{itemize}[nosep]
  \item Newton: \textbf{marginal} in 11/11 $(\beta, C)$ points (Earnshaw).
  \item TGP: \textbf{stable} in 11/11 points --- all physical modes
        $\omega^2 > 0$.
  \item Effect of $V_3$: mildly destabilizes 1 of 3 modes (toward marginal),
        but 2 modes remain stable.
  \item 5-gon ($N=5$): TGP creates a \emph{saddle} ($\omega^2 = -20.7$
        tangential, $+24.7$ radial) --- stabilization is
        \textbf{geometry-specific}, not universal.
\end{itemize}

\medskip
\noindent\textbf{Significance.}  TGP is the only known modification of
gravity that breaks Earnshaw for static $N$-body configurations
\emph{without rotation or external fields} --- purely from $V(\Phi)$.

% =====================================================================
\section{Result 2: Chaos suppression in the three-body problem}
\label{sec:chaos}
% =====================================================================

\begin{prediction}[Falsifiable]
For $\beta \gtrsim 0.025$ and initial conditions with close encounters,
TGP suppresses the maximal Lyapunov exponent by 35--45\%\textup{:}
\begin{equation}
  \frac{\lmax^{\rm TGP}}{\lmax^{\rm Newton}} \approx 0.55\text{--}0.65.
  \label{eq:ratio}
\end{equation}
\end{prediction}

\subsection{Methodology}

We set $G_{\rm Newton} = 4\pi$ so that the leading TGP term
$V_{\rm grad} = -4\pi C_i C_j / d$ reproduces Newton in the limit
$\beta \to 0$.  Kinetic energy is matched so that
$H_{\rm TGP} = H_{\rm Newton}$ on the same configuration (Burrau
pythagorean IC~\citep{Burrau1913,Szebehely1967} with $v = 0$; energy matching via velocity scaling or
random velocity generation with fixed RNG seed for reproducibility).

The maximal Lyapunov exponent $\lmax$ is computed via the Benettin
algorithm~\citep{Benettin1980}:
\begin{enumerate}[nosep]
  \item Symplectic leapfrog (Velocity Verlet) for the base trajectory.
  \item Coupled tangent map with the full analytical Jacobian
        $J_{i\alpha,j\beta} = \partial a_{i,\alpha}/\partial x_{j,\beta}$
        (chain rule through the Hessian of $I_Y$ in the distance variables).
  \item QR renormalization at regular intervals; accumulation of
        $\sum \log\|R_{kk}\|$.
\end{enumerate}

For the full spectrum ($k = 6N = 18$ exponents), Hamiltonian conservation
is verified: $|\sum_i \lambda_i| \sim 10^{-13}$.

\subsection{Beta scan with $V_2 + V_3$ (Burrau IC)}

\begin{center}
\begin{tabular}{ccccl}
  \toprule
  $\beta$ & ratio($V_2{+}V_3$) & ratio($V_2$) & $V_3$ effect & Classification \\
  \midrule
  0.020 & 1.09 & 1.91 & $-43\%$ & comparable \\
  0.050 & 0.52 & 0.67 & $-22\%$ & \textsc{suppress} \\
  0.080 & 0.58 & 0.91 & $-36\%$ & \textsc{suppress} \\
  0.120 & 0.63 & 0.99 & $-36\%$ & \textsc{suppress} \\
  0.200 & 0.65 & 0.85 & $-24\%$ & \textsc{suppress} \\
  0.350 & 0.60 & 0.76 & $-21\%$ & \textsc{suppress} \\
  \bottomrule
\end{tabular}
\end{center}

\begin{itemize}[nosep]
  \item 9/10 $\beta$-points: \textsc{suppress} (with $V_3$) vs 7/10
        ($V_2$-only).
  \item $\beta_{\rm crit}(V_3) \approx 0.025$ vs
        $\beta_{\rm crit}(V_2) \approx 0.042$ --- \textbf{$V_3$ lowers
        threshold by $\sim40\%$}.
  \item Convergence: ratio stable from $t = 2$ to $t = 6$.
\end{itemize}

\subsection{IC-dependence (multi-IC scan)}

\begin{center}
\begin{tabular}{lcc}
  \toprule
  \textbf{IC type} & \textbf{SUPPRESS fraction} & \textbf{Mechanism} \\
  \midrule
  Burrau (close encounters) & 75\% & Potential softening at small $d$ \\
  Equilateral (symmetric) & 0\% & No close encounters \\
  \bottomrule
\end{tabular}
\end{center}

The suppression mechanism is \emph{not} the repulsive barrier $\drep$
(which does not exist for typical $C$ values; it requires
$\beta \gg 1$).  Rather, TGP's higher-order terms ($1/d^2$, $1/d^3$)
\textbf{effectively soften the potential} at close encounters, reducing
the exponential divergence of nearby trajectories.  Symmetric ICs
(equilateral) never enter the close-encounter regime and thus show no
suppression.

\subsection{Caveats}

\begin{enumerate}[nosep]
  \item Finite-time Lyapunov exponents ($t_{\rm final} \leq 6$).
  \item Burrau IC with $v = 0$ requires random velocities for energy
        matching --- orbits depend on the RNG seed (seed 200 fixed for
        reproducibility).
  \item Leapfrog energy drift grows at long horizons; RK45 gives
        \texttt{success=False} at close encounters.
\end{enumerate}

% =====================================================================
\section{Result 3: Yukawa coupling from EFT projection of topological defect}
\label{sec:pathC}
% =====================================================================

\begin{result}[Path C closure]
The Yukawa source axiom (Path~B: $\delta_i = C_i\,e^{-\msp r}/(4\pi r)$)
is \textbf{derived} from the physics of the classical TGP defect via an EFT
projection.
\end{result}

\subsection{The problem}

The classical TGP field equation for a spherical
defect~\citep[cf.][for soliton physics]{Manton2004},
$g^2 g'' + g(g')^2 + (2/r)\,g^2\,g' = g^2(1-g)$, admits solutions
with boundary conditions $g(0) = g_0 < 1$, $g(\infty) = 1$.
Linearization at $g \to 1$ gives $(\nabla^2 + 1)\delta = 0$ where
$\delta = 1 - g$, yielding an \textbf{oscillatory} tail
$\delta(r) \sim \sin(r)/r$ (verified numerically: 7 sign changes).
This is \emph{not} Yukawa.

\subsection{Stabilization obstacle}

A naive vacuum-stabilization potential
$V_{\rm sb} = (\mu^2/2)(g-1)^2$ fails: the derivative
$V'_C(g) = (1-g)(g^2 - \mu^2)$ is negative for all $g \in (0,1)$
when $\mu^2 > 1$, destroying the defect's turning point.

\subsection{EFT resolution}

The screening mass $\msp$ enters not through a classical Lagrangian
modification, but through \textbf{loop corrections to the propagator}
(EFT mass gap).  The effective source coupling is the projection of the
classical defect profile onto the Yukawa Green's function:
\begin{equation}
  \boxed{
  \Ceff = \int_0^\infty \delta(r)\;e^{-\msp\,r}\;r\;\mathrm{d}r,
  }
  \label{eq:Ceff}
\end{equation}
where $\delta(r) = 1 - g(r)$ is the classical defect profile.

\subsection{Scaling laws}

Scanning over defect depth $g_0$ and coupling $\beta$:
\begin{equation}
  |\Ceff| \propto (1 - g_0)^{1.02},
  \qquad
  |\Ceff| \propto \beta^{-1.00}.
\end{equation}
The exponents are remarkably close to unity, giving a transparent physical
interpretation: deeper cores and weaker self-coupling produce larger
effective Yukawa charges.

\medskip\noindent
In the auxiliary stabilization picture used to motivate the EFT bridge, the
parameter $\mu^2$ satisfies
\[
  \mu^2 = 2(3\gamma - 2\beta) = 2\,\msp^2.
\]
It is therefore \emph{not} identical to Eq.~\eqref{eq:msp} squared; rather,
it is an auxiliary scale whose induced effective mass reproduces the physical
screening mass $\msp$ used by the $N$-body layer.

% =====================================================================
\section{Result 4: Classical soliton interaction is not Yukawa}
\label{sec:soliton-interaction}
% =====================================================================

\begin{result}[EFT mass gap required]
The classical overlap interaction of two TGP defects is oscillatory and
decays as a power law ($\sim d^{-1/2}$), not exponentially.  The Yukawa
force law used in $N$-body dynamics is a \textbf{quantum} (EFT) result,
not a classical one.
\end{result}

\subsection{Method}

Two defect profiles are placed at separation $d$ and their field overlap
is integrated on a 2D cylindrical grid $(\rho, z)$:
\begin{align}
  S(d) &= 2\pi \iint \delta_1\,\delta_2\;\rho\;\mathrm{d}\rho\,\mathrm{d}z
    & &\text{(overlap integral)},
    \label{eq:S}\\
  G(d) &= \iint \nabla\delta_1 \cdot \nabla\delta_2 \;\mathrm{d}^3x
    & &\text{(gradient cross integral)},
    \label{eq:G}\\
  V_{\rm cl}(d) &= G(d) + V''(1)\,S(d)
    & &\text{(classical interaction energy)}.
    \label{eq:Vcl}
\end{align}

\subsection{Results}

\begin{itemize}[nosep]
  \item $\delta(d)$ oscillates (7 sign changes) --- the classical defect
        tail is not Yukawa.
  \item $G(d)$ and $V''(1)\,S(d)$ cancel to $\sim93\%$ --- confirming
        the linearized equation $(\nabla^2 + 1)\delta = 0$.
  \item $V_{\rm cl}$ decays as $d^{-1/2}$ (power law), not $e^{-\msp d}$.
  \item $|V_{\rm cl}/V_Y|$ grows from $\sim16$ to $\sim10^8$ with distance
        --- completely different physics.
\end{itemize}

\medskip
\noindent\textbf{Significance.}  This closes the deductive chain:
\[
  g(r) \;\xrightarrow{\text{P5: EFT projection}}\; \Ceff
  \;\xrightarrow{\text{Path B}}\; V_Y(d)
  \;\xrightarrow{\text{P0: EOM}}\; \vF_i.
\]
The Yukawa law is justified \emph{only} through the EFT mass gap, not
by classical field overlap.

% =====================================================================
\section{Result 5: Semi-analytic multipole representation of $I_Y$}
\label{sec:multipole}
% =====================================================================

Using the Yukawa addition theorem, the double sum in the partial-wave
expansion of $I_Y$ collapses to a \textbf{single sum over $l$}:
\begin{equation}
  \boxed{
  I_Y = 4\pi \sum_{l=0}^{L_{\max}} (2l+1)\,P_l(\cos\omega)\,
    R_l(d_{13}, d_{23}, \msp),
  }
  \label{eq:IY-multipole}
\end{equation}
where $\omega$ is the angle at the vertex joining edges $d_{13}$ and $d_{23}$,
$P_l$ are Legendre polynomials, and $R_l$ is a 1D radial integral
computed by 3-segment Gauss--Legendre quadrature.  The key
simplification is the angular diagonalization:
$A_{ll'}(\omega) = \frac{4\pi}{2l+1}\,\delta_{ll'}\,P_l(\cos\omega)$.

\begin{center}
\begin{tabular}{ccc}
  \toprule
  $L_{\max}$ & Relative error (equilateral) & Relative error (scalene) \\
  \midrule
  3  & 0.23\% & 2.1\% \\
  10 & 0.12\% & 0.22\% \\
  15 & 0.006\% & 0.04\% \\
  \bottomrule
\end{tabular}
\end{center}

The Feynman 2D integral (Sec.~\ref{sec:eom}) remains faster per evaluation
with caching, but the multipole form is tabulable and analytically
transparent.

% =====================================================================
\section{Result 6: $V_3/V_2$ regime map}
\label{sec:v3v2}
% =====================================================================

The dimensionless parameter $t = \msp\,d$ controls the importance of
three-body forces relative to pairwise interactions.

\begin{center}
\begin{tabular}{lcl}
  \toprule
  \textbf{Regime} & $|V_3/V_2|$ & \textbf{Description} \\
  \midrule
  Coulomb ($t < 0.5$) & $\sim 1.3\%$ & Saturates (constant); $V_3$ perturbative \\
  Transition ($t \sim 1$) & $\sim 14\%$ & Peak; $V_3$ significant \\
  Yukawa ($t > 3$) & $< 0.2\%$ & Exponentially screened \\
  \bottomrule
\end{tabular}
\end{center}

\noindent Additional findings:
\begin{itemize}[nosep]
  \item \textbf{Lyapunov regime} ($\beta = 0.02$--$0.35$,
        $\msp d \sim 0.3$--$1.0$): $|V_3/V_2| = 0.8$--$1.0\%$ ---
        perturbative but sufficient to lower $\beta_{\rm crit}$ by $40\%$.
  \item \textbf{Geometry}: compact configurations ($d/2$) have $7\times$
        larger $V_3/V_2$ than extended ($2d$) --- close encounters amplify
        three-body effects.
  \item \textbf{$C$-scaling}: $V_3/V_2 \propto C$ (linear, verified to
        $<2\%$ deviation) --- consistent with $V_3 \propto C^3$ vs
        $V_2 \propto C^2$.
  \item \textbf{Physical TGP} ($\gamma \sim \Lambda_{\rm eff} \sim 10^{-52}$):
        $\msp d \ll 1$ --- Coulomb regime applies for all macroscopic
        separations.
\end{itemize}

% =====================================================================
\section{Result 7: Analytical equilibria, normal modes, and Hill criterion}
\label{sec:analytical}
% =====================================================================

\subsection{Exact equilibrium distances}

Setting $V_2'(d) = 0$ for equal-mass bodies ($C_1 = C_2 = C$) yields a
quadratic in~$d$:
\begin{equation}
  \boxed{
  d^2 - 4\beta\,d + 18\gamma\,C = 0.
  }
  \label{eq:equilibrium-condition}
\end{equation}
The two roots give the \emph{repulsive barrier} and the \emph{confining
well}:
\begin{equation}
  \drep = 2\beta - \sqrt{4\beta^2 - 18\gamma C},
  \qquad
  \dwell = 2\beta + \sqrt{4\beta^2 - 18\gamma C},
  \label{eq:drep-dwell}
\end{equation}
existing whenever $\beta > \sqrt{9\gamma C/2}$.  For unequal masses, the
condition generalizes to $d^2 - 4\beta\,d + 9\gamma(C_1+C_2) = 0$.
At threshold ($4\beta^2 = 18\gamma C$), the two equilibria merge:
$\drep = \dwell = 2\beta$.  Far from threshold: $\drep \to 0$ and
$\dwell \to 4\beta$.  Verified numerically: analytical = numerical to
machine precision.

\subsection{Closed-form normal-mode frequency}

The radial stiffness at $\dwell$ is
\begin{equation}
  V_2''(d) = \frac{8\pi C^2}{d^3}
    \biggl[-1 + \frac{6\beta}{d} - \frac{18\gamma C}{d^2}\biggr].
  \label{eq:V2pp}
\end{equation}
For an isolated pair with reduced mass $\mu = C/2$, the radial breathing
frequency is
\begin{equation}
  \boxed{
  \omega_{\rm rad}^2 = \frac{2\,V_2''(\dwell)}{C}
  = \frac{16\pi C}{\dwell^3}
    \biggl[-1 + \frac{6\beta}{\dwell} - \frac{18\gamma C}{\dwell^2}\biggr].
  }
  \label{eq:omega-rad}
\end{equation}
Verified against finite-difference Hessian: relative error $<10^{-6}$.

\subsection{TGP Hill sphere}

The classical Hill sphere~\citep{Murray1999,Hamilton1992} is the region
around a body in which its gravity dominates over the tidal field of
a more massive companion.
The tidal field of a central body $C_{\rm cen}$ at distance~$d$ in TGP is
\begin{equation}
  \Omega_{\rm tidal}^2(d)
  = \frac{8\pi C_{\rm cen}}{d^3}
    \biggl[1 - \frac{6\beta}{d} + \frac{18\gamma\,C_{\rm cen}}{d^2}\biggr].
  \label{eq:tidal}
\end{equation}
The $\beta$-repulsion \emph{weakens} the tidal field (negative correction to
the bracket), while $\gamma$-confinement \emph{strengthens} it.  The Hill
radius follows from balancing the tidal acceleration against self-gravity
of a body of mass $C_{\rm body}$:
\begin{equation}
  \boxed{
  R_H^3 = \frac{4\pi\,C_{\rm body}}{|\Omega_{\rm tidal}^2(d)|}.
  }
  \label{eq:hill}
\end{equation}
Scanning $\beta \in [0.01, 2.0]$ at $d = 5$, $C_{\rm body} = 0.1$,
$C_{\rm cen} = 1.0$:

\begin{center}
\begin{tabular}{ccl}
  \toprule
  $\beta$ & $R_H^{\rm TGP}/R_H^{\rm Newton}$ & Effect \\
  \midrule
  0.01 & 1.002 & $\sim$same \\
  0.10 & 1.017 & slightly larger \\
  0.50 & 1.096 & 10\% larger \\
  1.00 & 1.244 & 24\% larger \\
  2.00 & 2.924 & $3\times$ larger \\
  \bottomrule
\end{tabular}
\end{center}

\noindent\textbf{Result:} TGP generically \textbf{expands} the Hill sphere
($R_H^{\rm TGP} > R_H^{\rm Newton}$), making satellite systems more stable
against tidal disruption.

\subsection{Perturbative $V_3$ correction to equilibria}

The three-body force shifts $\dwell$ by
\begin{equation}
  \delta d = -\frac{F_3(\dwell)}{3\,V_2''(\dwell)},
  \label{eq:delta-d}
\end{equation}
where $F_3 = -\partial V_3/\partial d|_{\rm breathing}$.  Numerically:
$|\delta d/\dwell| < 0.2\%$ for all tested $(\beta, C)$ --- the
correction is perturbative.  The sign is positive: \textbf{$V_3$ shifts
$\dwell$ outward} (slightly larger equilibrium distance).

% =====================================================================
\section{Supplementary result: Closed orbits and Lagrange stability}
\label{sec:orbits}
% =====================================================================

Two classical periodic configurations were tested under TGP perturbation:

\begin{enumerate}[nosep]
  \item \textbf{Figure-8 choreography} (Chenciner--Montgomery):
        survives TGP to $\beta = 0.10$ ($C = 0.3$).  Period shift:
        Newton $T = 3.115$, TGP $T = 3.283$ (+5.4\%).  Deformation
        scales as $\Delta r_{\max} \sim \beta^{0.65}$ (perturbative).
        With $V_3$: $V_3/V_2 = 0.82\%$.
  \item \textbf{Lagrange point $L_4$}: position shift $\sim 10^{-5}$;
        stability preserved ($q_{\rm crit}$ unchanged); Poincar\'e
        sections show regular orbits.  Effective-potential curvature
        changes by $< 10^{-4}$.
\end{enumerate}

\noindent\textbf{Conclusion.}  Periodic orbits are \emph{robust} under TGP
corrections --- the perturbation is smooth and non-destructive.

% =====================================================================
\section{Result 8: Unequal masses, phase diagram, escape velocity, breathing mode}
\label{sec:unequal}
% =====================================================================

\begin{result}[Unequal-mass confinement and phase diagram]
For $C_1\neq C_2$ the force equation $F(d)=0$ yields
\begin{equation}\label{eq:unequal}
  d^{2} - 4\beta\,d + 9\gamma(C_1{+}C_2) = 0,
\end{equation}
with roots $d_{\rm rep,well} = 2\beta \mp \sqrt{4\beta^{2}-9\gamma(C_1{+}C_2)}$.
\end{result}

\noindent Eq.~\eqref{eq:unequal} reduces to eq.~\eqref{eq:equilibrium-condition}
for $C_1{=}C_2{=}C$.

\subsection{Phase diagram}

Confinement exists if and only if the discriminant is positive:
\begin{equation}\label{eq:phase}
\boxed{\;
  \beta > \beta_{\rm crit} \;=\; \tfrac{3}{2}\,\sqrt{\gamma\,(C_1{+}C_2)}\,,
  \qquad\text{equivalently}\quad
  C_1{+}C_2 < \tfrac{4}{9}\,\frac{\beta^{2}}{\gamma}\,.
\;}
\end{equation}
At $\beta=\beta_{\rm crit}$ the two roots coalesce:
$\drep = \dwell = 2\beta_{\rm crit}$ (saddle-node bifurcation).
For fixed $C_1$, the maximum mass ratio for which equilibria persist is
\begin{equation}
  q_{\max} = \frac{C_{2,\max}}{C_1}
  = \frac{4\beta^{2}}{9\gamma\,C_1} - 1.
\end{equation}
Numerical verification (\texttt{ex211}, grid $10\times10$ in
$(\beta, C_{\rm sum})$): 100\% agreement with analytical boundary.

\subsection{Escape velocity}

The binding energy at $\dwell$ is $E_{\rm bind} = |V_2(\dwell)|$,
so the escape velocity in the reduced-mass frame is
\begin{equation}\label{eq:vesc}
\boxed{\;
  v_{\rm esc} = \sqrt{\frac{2\,|V_2(\dwell)|}{\mu}}\,,
  \qquad \mu = \frac{C_1 C_2}{C_1{+}C_2}\,.
\;}
\end{equation}
Substituting $V_2(\dwell) = -A/\dwell + B/\dwell^{2} - G/\dwell^{3}$
gives a fully closed-form expression in $(\beta,\gamma,C_1,C_2)$.
Scaling: $v_{\rm esc}$ decreases monotonically with $\beta$
(deeper well at smaller $\beta$, but $\dwell$ also shrinks;
net effect: $v_{\rm esc} \propto \beta^{-0.5}$ approximately).

\subsection{Three-body breathing mode}

For three equal masses ($C_1{=}C_2{=}C_3{=}C$) in an equilateral triangle
with side~$d$, the effective potential for the breathing coordinate is
\begin{equation}
  U_{\rm eff}(d) \;=\; 3\,V_2(d) \;+\; V_3(d,d,d).
\end{equation}
The reduced mass for symmetric breathing (all distances change equally) is
$\mu_{\rm br} = C$, giving
\begin{equation}\label{eq:omega-br}
  \omega^{2}_{\rm br} = \frac{U_{\rm eff}''(\dwell)}{C}
  = \frac{3\,V_2''(\dwell) + V_3''(\dwell)}{C}\,.
\end{equation}
Without $V_3$, this simplifies to
$\omega^{2}_{\rm br} = 3\,V_2''(\dwell)/C = \tfrac{3}{2}\,\omega^{2}_{\rm rad}$
(the breathing mode is \textbf{50\% stiffer} than the two-body radial mode;
physically, three pairs contribute to the restoring force while the breathing
reduced mass is $C$, not $C/2$).

The $V_3$ correction to $\omega^{2}_{\rm br}$ is \textbf{$-1.3\%$}
(\texttt{ex211}, $\beta{=}\gamma{=}1$, $C{=}0.1$) --- perturbative
and consistent with the $V_3/V_2$ regime analysis (Result~6).

\medskip
\noindent\textbf{Inward barrier.}
The energy barrier for the breathing mode to reach $\drep$ from $\dwell$ is
$\Delta U = U_{\rm eff}(\drep) - U_{\rm eff}(\dwell)$.
Numerically: $\Delta U / |U_{\rm eff}(\dwell)| \approx 9.4$
for $\beta{=}1$, $C{=}0.1$ --- the inward barrier is roughly ten times
the binding energy, confirming that the repulsive wall at $\drep$ is
highly effective.

% =====================================================================
\section{Summary and deductive chain}
\label{sec:summary}
% =====================================================================

The complete deductive chain from the TGP field Lagrangian to
$N$-body dynamics is:

\begin{center}
\begin{tabular}{rl}
  \textbf{Layer} & \textbf{Content} \\
  \midrule
  Field equation & $K_{\rm sub}(g)\;\Box g + \ldots = V'(g)$
    \quad ($K_{\rm sub} = g^2$) \\
  Defect profile & $g(r)$: ODE solution, oscillatory tail $\sin(r)/r$ \\
  EFT projection (P5) & $\Ceff = \int \delta\,e^{-\msp r}\,r\,\mathrm{d}r$ \\
  Yukawa source (Path B) & $\delta_i = \Ceff\,e^{-\msp r}/(4\pi r)$ \\
  $V_2$ (pairwise) & Eq.~\eqref{eq:V2}: gradient + repulsion + confinement \\
  $V_3$ (three-body) & Eq.~\eqref{eq:V3}: Feynman integral $I_Y$ \\
  $N$-body EOM & $C_i\ddot\vx_i = -\nabla_{\vx_i}(V_2 + V_3)$ \\
  \midrule
  \multicolumn{2}{c}{\textit{Novel results from this chain:}} \\
  \midrule
  Result 1 & Earnshaw violation (Sec.~\ref{sec:earnshaw}) \\
  Result 2 & Chaos suppression 35--45\% (Sec.~\ref{sec:chaos}) \\
  Result 3 & $\Ceff$ from EFT projection (Sec.~\ref{sec:pathC}) \\
  Result 4 & Classical $\neq$ Yukawa (Sec.~\ref{sec:soliton-interaction}) \\
  Result 5 & $I_Y$ multipole form (Sec.~\ref{sec:multipole}) \\
  Result 6 & $V_3/V_2$ regime map (Sec.~\ref{sec:v3v2}) \\
  Result 7 & Analytical equilibria + Hill sphere (Sec.~\ref{sec:analytical}) \\
  Result 8 & Phase diagram, $v_{\rm esc}$, breathing mode (Sec.~\ref{sec:unequal}) \\
\end{tabular}
\end{center}

\bigskip

All results are supported by a regression suite of 59 numerical tests
(102 active example scripts) passing consistently.  The implementation
is available in the \texttt{nbody/} package of TGP~v1.

\paragraph{Context.}
The pairwise potential $V_2$ connects to the three regimes of TGP
(gradient, repulsive, confining) identified in the parent
theory~\citep{Serafin2026}.  The $N$-body dynamics presented here are
\emph{exact} consequences of the field Lagrangian (no tunable
parameters beyond $\beta$, $\gamma$, $C_i$), distinguishing TGP from
phenomenological modified-gravity $N$-body
approaches~\citep{Puchwein2013,Li2012,Winther2015}.  The Earnshaw
violation (Result~1) and chaos suppression (Result~2) are qualitatively
new predictions with no counterpart in Newtonian gravity; the EFT
derivation of $\Ceff$ (Result~3) bridges the gap between the
topological defect (UV) and the Yukawa source (IR).

\paragraph{Outlook.}
Future work includes: (i)~extending to relativistic corrections using
the TGP metric ($\gamma_{\rm PPN} = \beta_{\rm PPN} = 1$, per
\citealt{Will2014}); (ii)~cosmological $N$-body simulations
incorporating scale-dependent $G_{\rm eff}(k,a)$;
(iii)~exploring the Efimov-like tower that may arise from
$V_3$ near the two-body resonance at $\beta_{\rm crit}$
\citep{Efimov1970,Braaten2006};
(iv)~applying the analytical equilibria (Result~7) and Hill criterion
to exoplanetary stability~\citep{Murray1999,Hamilton1992}.

\paragraph{Reproducibility.}
All numerical results are reproduced by
\texttt{python reproduce\_all.py -{}-quick}
(10/10 stages PASS in ${\sim}100$\,s on a standard laptop).
See \texttt{README.md} for the full script manifest
and the supplementary material for data files.

\bibliography{tgp_nbody}

\end{document}
